\section{Results}\label{results}

\begin{quote}
\textbf{Note:} This file documents the result structure and expected
findings. Populate with actual values by running:

\begin{Shaded}
\begin{Highlighting}[]
\ExtensionTok{python}\NormalTok{ experiments/run\_attacks.py }\AttributeTok{{-}{-}epsilon}\NormalTok{ 3.0 }\AttributeTok{{-}{-}n{-}docs}\NormalTok{ 50}
\ExtensionTok{python}\NormalTok{ experiments/ablations.py }\AttributeTok{{-}{-}n{-}docs}\NormalTok{ 20}
\ExtensionTok{python}\NormalTok{ experiments/calibrate\_epsilon.py}
\end{Highlighting}
\end{Shaded}

All claims must be traceable to a file in \texttt{experiments/results/}.
\end{quote}

\subsection{1. Privacy-Utility Curve}\label{privacy-utility-curve}

Produced by \texttt{experiments/calibrate\_epsilon.py}. For each ε ∈
\{0.5, 1.0, 2.0, 3.0, 5.0, 10.0\}:

{\def\LTcaptype{none} % do not increment counter
\begin{longtable}[]{@{}
  >{\raggedright\arraybackslash}p{(\linewidth - 10\tabcolsep) * \real{0.0333}}
  >{\raggedright\arraybackslash}p{(\linewidth - 10\tabcolsep) * \real{0.1444}}
  >{\raggedright\arraybackslash}p{(\linewidth - 10\tabcolsep) * \real{0.1444}}
  >{\raggedright\arraybackslash}p{(\linewidth - 10\tabcolsep) * \real{0.1667}}
  >{\raggedright\arraybackslash}p{(\linewidth - 10\tabcolsep) * \real{0.2667}}
  >{\raggedright\arraybackslash}p{(\linewidth - 10\tabcolsep) * \real{0.2444}}@{}}
\toprule\noalign{}
\begin{minipage}[b]{\linewidth}\raggedright
ε
\end{minipage} & \begin{minipage}[b]{\linewidth}\raggedright
σ (Gaussian)
\end{minipage} & \begin{minipage}[b]{\linewidth}\raggedright
Inversion F1
\end{minipage} & \begin{minipage}[b]{\linewidth}\raggedright
Utility Ratio
\end{minipage} & \begin{minipage}[b]{\linewidth}\raggedright
Passes H₁ (F1 ≤ 0.09)
\end{minipage} & \begin{minipage}[b]{\linewidth}\raggedright
Passes H₂ (U ≥ 0.85)
\end{minipage} \\
\midrule\noalign{}
\endhead
\bottomrule\noalign{}
\endlastfoot
0.5 & \textasciitilde6.60 & \emph{TBD} & \emph{TBD} & \emph{TBD} &
\emph{TBD} \\
1.0 & \textasciitilde3.30 & \emph{TBD} & \emph{TBD} & \emph{TBD} &
\emph{TBD} \\
2.0 & \textasciitilde1.65 & \emph{TBD} & \emph{TBD} & \emph{TBD} &
\emph{TBD} \\
\textbf{3.0} & \textbf{\textasciitilde1.10} & \textbf{\emph{TBD}} &
\textbf{\emph{TBD}} & \textbf{\emph{TBD}} & \textbf{\emph{TBD}} \\
5.0 & \textasciitilde0.66 & \emph{TBD} & \emph{TBD} & \emph{TBD} &
\emph{TBD} \\
10.0 & \textasciitilde0.33 & \emph{TBD} & \emph{TBD} & \emph{TBD} &
\emph{TBD} \\
\end{longtable}
}

σ values computed analytically: σ = √(2 ln(1.25/1e-5)) / ε ≈ 3.298 / ε.

Result source: \texttt{experiments/results/calibration.json}

\subsection{2. Full Attack Battery at ε =
3.0}\label{full-attack-battery-at-ux3b5-3.0}

Result source: \texttt{experiments/results/attacks\_eps3.0.json}

\subsubsection{Attack 1 --- Verbatim Leak}\label{attack-1-verbatim-leak}

{\def\LTcaptype{none} % do not increment counter
\begin{longtable}[]{@{}lll@{}}
\toprule\noalign{}
Category & Literal Leak Rate & Fuzzy Leak Rate (threshold 85) \\
\midrule\noalign{}
\endhead
\bottomrule\noalign{}
\endlastfoot
NAME & \emph{TBD} & \emph{TBD} \\
DATE & \emph{TBD} & \emph{TBD} \\
SSN & \emph{TBD} & \emph{TBD} \\
COMPOUND\_CODE & \emph{TBD} & \emph{TBD} \\
SITE\_ID & \emph{TBD} & \emph{TBD} \\
AE\_GRADE & \emph{TBD} & \emph{TBD} \\
EFFICACY\_VALUE & \emph{TBD} & \emph{TBD} \\
\textbf{Overall} & \textbf{\emph{TBD}} & \textbf{\emph{TBD}} \\
\end{longtable}
}

\textbf{Expected finding:} Safe Harbor stripping should drive literal
leak rates to near zero for structured identifiers (NAME, DATE, SSN,
SITE\_ID). Compound codes and quasi-identifiers are handled by NGSP;
their residual rates test the neural components' marginal value.

\subsubsection{Attack 2 --- Cross-Encoder
Similarity}\label{attack-2-cross-encoder-similarity}

{\def\LTcaptype{none} % do not increment counter
\begin{longtable}[]{@{}ll@{}}
\toprule\noalign{}
Metric & Value \\
\midrule\noalign{}
\endhead
\bottomrule\noalign{}
\endlastfoot
mean\_sim & \emph{TBD} \\
median\_sim & \emph{TBD} \\
p95\_sim & \emph{TBD} \\
fraction\_above\_threshold (0.85) & \emph{TBD} \\
\end{longtable}
}

\textbf{Expected finding:} The abstract-extractable path (query
synthesis) should produce low similarity scores because the synthesized
query captures task intent but drops entity specifics. The DP path may
produce higher similarity scores since the proxy text retains sentence
structure.

\subsubsection{Attack 3 --- Trained Inversion (Primary
Result)}\label{attack-3-trained-inversion-primary-result}

{\def\LTcaptype{none} % do not increment counter
\begin{longtable}[]{@{}ll@{}}
\toprule\noalign{}
Metric & Value \\
\midrule\noalign{}
\endhead
\bottomrule\noalign{}
\endlastfoot
overall\_f1 & \emph{TBD} \\
baseline\_random\_f1 & \emph{TBD} \\
control\_shuffle\_f1 & \emph{TBD} \\
n\_train & \emph{TBD} \\
n\_eval & \emph{TBD} \\
\end{longtable}
}

\textbf{H₁ verdict:} \emph{TBD} (threshold: overall\_f1 ≤ 0.09)

\textbf{Expected finding:} The shuffle-pair control F1 should be close
to the random baseline, confirming that any learned signal above
baseline is real and not a training artifact. If the main model F1 is
also near baseline, NGSP successfully breaks the proxy→span mapping.

\subsubsection{Attack 4 --- Membership
Inference}\label{attack-4-membership-inference}

{\def\LTcaptype{none} % do not increment counter
\begin{longtable}[]{@{}llll@{}}
\toprule\noalign{}
Entity & AUC & n\_positive & n\_negative \\
\midrule\noalign{}
\endhead
\bottomrule\noalign{}
\endlastfoot
\emph{top entity 1} & \emph{TBD} & \emph{TBD} & \emph{TBD} \\
\emph{top entity 2} & \emph{TBD} & \emph{TBD} & \emph{TBD} \\
\emph{top entity 3} & \emph{TBD} & \emph{TBD} & \emph{TBD} \\
\textbf{mean\_auc} & \textbf{\emph{TBD}} & & \\
\end{longtable}
}

\textbf{Expected finding:} AUC near 0.5 indicates the proxy text does
not expose whether a specific entity was mentioned in the original
document. AUC significantly above 0.5 would indicate residual membership
signal.

\subsubsection{Attack 5 --- Downstream
Utility}\label{attack-5-downstream-utility}

{\def\LTcaptype{none} % do not increment counter
\begin{longtable}[]{@{}ll@{}}
\toprule\noalign{}
Metric & Value \\
\midrule\noalign{}
\endhead
\bottomrule\noalign{}
\endlastfoot
mean\_original\_score & \emph{TBD} \\
mean\_proxy\_score & \emph{TBD} \\
utility\_ratio & \emph{TBD} \\
n\_docs\_evaluated & \emph{TBD} \\
\end{longtable}
}

\textbf{H₂ verdict:} \emph{TBD} (threshold: utility\_ratio ≥ 0.85)

\subsection{3. Ablation Study}\label{ablation-study}

Result source: \texttt{experiments/results/ablations.json}

{\def\LTcaptype{none} % do not increment counter
\begin{longtable}[]{@{}
  >{\raggedright\arraybackslash}p{(\linewidth - 6\tabcolsep) * \real{0.2000}}
  >{\raggedright\arraybackslash}p{(\linewidth - 6\tabcolsep) * \real{0.3143}}
  >{\raggedright\arraybackslash}p{(\linewidth - 6\tabcolsep) * \real{0.2857}}
  >{\raggedright\arraybackslash}p{(\linewidth - 6\tabcolsep) * \real{0.2000}}@{}}
\toprule\noalign{}
\begin{minipage}[b]{\linewidth}\raggedright
Configuration
\end{minipage} & \begin{minipage}[b]{\linewidth}\raggedright
Verbatim Literal Leak
\end{minipage} & \begin{minipage}[b]{\linewidth}\raggedright
Verbatim Fuzzy Leak
\end{minipage} & \begin{minipage}[b]{\linewidth}\raggedright
Inversion F1
\end{minipage} \\
\midrule\noalign{}
\endhead
\bottomrule\noalign{}
\endlastfoot
Safe Harbor only & \emph{TBD} & \emph{TBD} & \emph{TBD} \\
Safe Harbor + synthesis & \emph{TBD} & \emph{TBD} & \emph{TBD} \\
Safe Harbor + DP only & \emph{TBD} & \emph{TBD} & \emph{TBD} \\
Full system & \emph{TBD} & \emph{TBD} & \emph{TBD} \\
\end{longtable}
}

\textbf{Expected finding (H₃):} No single component should achieve both
privacy and utility thresholds simultaneously. Safe Harbor alone
eliminates structured PHI but leaves quasi-identifiers exposed. Query
synthesis breaks the inversion signal for abstract-extractable inputs
but cannot handle all document types. DP alone provides formal
guarantees but may degrade utility at tight ε values. The full system
combines all three.

\textbf{Negative result to report honestly:} If the ablation shows that
Safe Harbor alone already achieves F1 ≤ 0.09 (because the synthetic
corpus's sensitive spans are dominated by Safe Harbor categories), that
is a negative result for H₃ and should be stated clearly. It would mean
NGSP provides marginal privacy improvement over simpler baselines on
this corpus --- a finding worth reporting.

\subsection{4. Comment Audit}\label{comment-audit}

Result: \texttt{paper/code\_audit.md}

Running \texttt{python3\ scripts/check\_comments.py\ -\/-path\ src/} as
of Phase 3 completion:

\begin{verbatim}
OK — all defs have a one-line comment directly above. (src)
\end{verbatim}

Zero violations across all 26 Python files in \texttt{src/}.
